\newpage
\section{Phương pháp đề xuất}

\subsection{Pipeline tổng quát của hệ thống}

Hệ thống dehaze được xây dựng theo một pipeline tuần tự gồm các bước chính sau: chuẩn hóa ảnh đầu vào, tính dark channel, ước lượng atmospheric light, dựng transmission thô từ ảnh đã chuẩn hóa theo \(A\), tinh chỉnh transmission bằng guided filter, sau đó khôi phục ảnh cuối cùng. Về mặt học thuật, đây vẫn là pipeline DCP cổ điển; điểm khác biệt chính nằm ở cách tổ chức tính toán để phù hợp với plugin và ở khả năng thay \(\omega\) cố định bằng \(\omega(x)\) thích nghi.

Ảnh khôi phục \(J(x)\) được tính theo công thức:
\begin{equation}
J(x)=\frac{I(x)-A}{\max(t(x),t_0)}+A,
\end{equation}
trong đó \(t_0\) là ngưỡng transmission tối thiểu để tránh chia cho giá trị quá nhỏ. Trong hiện thực hiện tại, ngưỡng này được đặt là \(t_0 = 0.1\).

\begin{figure}[H]
  \centering
  \begin{tikzpicture}[
    box/.style={draw, rounded corners=2pt, align=center, minimum width=2.45cm, minimum height=0.95cm, fill=blue!7},
    aux/.style={draw, rounded corners=2pt, align=center, minimum width=2.2cm, minimum height=0.85cm, fill=green!8},
    arrow/.style={->, thick}
  ]
    \node[box] (input) at (0,0) {Ảnh hazy\\đầu vào};
    \node[box] (dark) at (3.1,0) {Dark channel};
    \node[box] (airlight) at (6.5,0) {Ước lượng\\atmospheric light \(A\)};
    \node[box] (slope) at (10.4,0) {Transmission thô\\và guided filter};
    \node[box] (recover) at (14.6,0) {Recovery\\ảnh đầu ra};

    \node[aux] (adaptive) at (10.4,-2.45) {Adaptive omega\\(tùy chọn)};
    \node[aux] (strength) at (14.6,-2.45) {Strength};

    \draw[arrow] (input) -- (dark);
    \draw[arrow] (dark) -- (airlight);
    \draw[arrow] (airlight) -- (slope);
    \draw[arrow] (slope) -- node[above=3pt] {\(t(x)\)} (recover);
    \draw[arrow] (adaptive.north) -- node[left=2pt] {\(\omega(x)\)} (slope.south);
    \draw[arrow] (strength.north) -- (recover.south);
  \end{tikzpicture}
  \caption{Sơ đồ pipeline tổng quát của hệ thống dehaze. Nhánh adaptive omega chỉ được kích hoạt khi người dùng bật chế độ tương ứng}
\end{figure}

\subsection{Hiện thực dark channel và atmospheric light}

Trong mã nguồn, ảnh đầu vào được đưa về miền \([0,1]\) theo từng kênh màu. Với mỗi điểm ảnh, hệ thống tạo một tín hiệu trung gian
\[
\texttt{darkInput}(x) = \min(r(x), g(x), b(x)),
\]
tương ứng với bước lấy cực tiểu theo kênh trong định nghĩa dark channel. Song song với đó, một ảnh hướng dẫn mức xám cũng được xây dựng từ độ chói:
\[
\texttt{guidance}(x)=0.299\,r(x)+0.587\,g(x)+0.114\,b(x).
\]
Ảnh xám này không dùng để tạo dark channel mà được giữ lại cho bước guided filter phía sau.

Dark channel được tạo bằng hàm \texttt{MinFilter} hai chiều với bán kính \(7\), tương ứng patch kích thước \(15\times15\). Về mặt học thuật, thao tác này chính là xấp xỉ trực tiếp
\[
\min_{y \in \Omega(x)} \left(\min_c I^c(y)\right).
\]
Cách cài đặt trong plugin không dùng thư viện xử lý ảnh ngoài mà hiện thực phép lọc cực tiểu theo hai lượt ngang và dọc, đủ đơn giản để kiểm soát trong môi trường plugin.

Sau khi có dark channel, atmospheric light được ước lượng bằng một quy trình bám sát ý tưởng của DCP. Trước hết, hệ thống tạo histogram 256 mức của dark channel và lấy ngưỡng tương ứng với nhóm khoảng \(0.1\%\) điểm ảnh sáng nhất trong dark channel. Sau đó, trong tập ứng viên vượt ngưỡng này, thuật toán chọn điểm có tổng cường độ \(R+G+B\) lớn nhất trên ảnh gốc. Giá trị RGB tại điểm này được dùng làm \(A\). Ngoài ra, mỗi thành phần của \(A\) còn được chặn dưới bởi \(0.1\) để tránh trường hợp mẫu số quá nhỏ khi chuẩn hóa ảnh theo atmospheric light.

\subsection{Xây dựng transmission map}

Khi \(A\) đã được xác định, ảnh đầu vào được chuẩn hóa theo atmospheric light trên từng kênh:
\begin{equation}
\frac{I^c(x)}{A^c}.
\end{equation}
Trong hiện thực, plugin không ghép \(\omega\) vào ngay tại bước này. Thay vào đó, nó tạo một đại lượng trung gian:
\[
\texttt{rawTransmissionInput}(x) = \min_c \frac{I^c(x)}{A^c},
\]
rồi tiếp tục áp dụng \texttt{MinFilter} để thu được bản đồ
\[
\texttt{normalizedDarkChannel}(x) \approx \min_{y\in\Omega(x)}\left(\min_c \frac{I^c(y)}{A^c}\right).
\]
Đại lượng này chính là phần nhân với \(\omega\) trong công thức transmission thô của DCP.

Tiếp theo, \texttt{normalizedDarkChannel} được tinh chỉnh bằng guided filter với ảnh hướng dẫn là \texttt{guidance}. Kết quả của bước này được mã nguồn lưu dưới tên \texttt{transmissionSlope}. Tên gọi này phản ánh đúng vai trò của nó trong hiện thực: đây chưa phải transmission cuối cùng, mà là phần hệ số sẽ được nhân với \(\omega\) hoặc \(\omega(x)\) về sau. Cách tách riêng như vậy cho phép thay đổi mức dehaze bằng thanh \texttt{Strength} mà không phải xây lại toàn bộ transmission từ đầu.

Đoạn mã dưới đây thể hiện cách xây dựng \texttt{transmissionSlope} từ ảnh đã chuẩn hóa theo atmospheric light:

\begin{lstlisting}[style=csharpblock,caption={Trích đoạn dựng transmissionSlope từ ảnh đã chuẩn hóa theo atmospheric light}]
for (int i = 0; i < length; i++)
{
    ColorBgra32 color = source[i];
    float r = (color.R / 255f) / atmosphericLight.R;
    float g = (color.G / 255f) / atmosphericLight.G;
    float b = (color.B / 255f) / atmosphericLight.B;

    rawTransmissionInput[i] = MathF.Min(r, MathF.Min(g, b));
}

float[] normalizedDarkChannel =
    MinFilter(rawTransmissionInput, width, height, DarkChannelRadius);

return GuidedFilter(
    guidance,
    normalizedDarkChannel,
    width,
    height,
    GuidedFilterRadius,
    GuidedFilterEpsilon);
\end{lstlisting}

\subsection{Khôi phục ảnh và điều khiển mức dehaze}

Ảnh cuối cùng được khôi phục trong hàm \texttt{RecoverPixel}. Với mỗi điểm ảnh, plugin tính
\begin{equation}
t(x)=\mathrm{clamp}\left(1-\omega_{\text{eff}}(x)\,\texttt{transmissionSlope}(x),\; 0.1,\; 1.0\right),
\end{equation}
trong đó \(\omega_{\text{eff}}(x)\) là hệ số dehaze thực sự được dùng tại điểm ảnh.

Nếu chọn chế độ cố định, \(\omega_{\text{eff}}(x)\) là một hằng số:
\[
\omega_{\text{eff}} = 0.95 \times \texttt{strengthFactor},
\]
với \(\texttt{strengthFactor} = \frac{\texttt{Strength}}{100}\). Nếu bật chế độ adaptive, hệ thống dùng:
\[
\omega_{\text{eff}}(x)=\omega_{\mathrm{map}}(x)\times \texttt{strengthFactor}.
\]
Ở đây, \texttt{Strength} không chỉnh trực tiếp màu của ảnh đầu ra. Tham số này chỉ thay đổi mức loại bỏ haze trước bước recovery thông qua \(\omega_{\text{eff}}\). Vì vậy, thanh trượt thực chất điều khiển mức can thiệp vào transmission, thay vì đóng vai trò như một bộ trộn hậu kỳ giữa ảnh gốc và ảnh đã xử lý.

Sau khi có \(t(x)\), từng kênh màu được phục hồi theo đúng dạng của mô hình khí quyển:
\[
J^c(x)=\frac{I^c(x)-A^c}{t(x)} + A^c,
\]
rồi được chặn lại về \([0,1]\). Cách viết trong mã nguồn theo từng kênh giúp phép tính bám sát công thức gốc nhưng vẫn đơn giản để tối ưu và tái sử dụng cho preview lẫn render toàn ảnh.

Đoạn mã sau thể hiện cách phục hồi một điểm ảnh trong hệ thống:

\begin{lstlisting}[style=csharpblock,caption={Trích đoạn RecoverPixel và cách Strength đi vào recovery}]
float transmission =
    Math.Clamp(1f - (effectiveOmega * transmissionSlope), 0.10f, 1f);

float r = RecoverChannel(sourcePixel.R / 255f, atmosphericLight.R, transmission);
float g = RecoverChannel(sourcePixel.G / 255f, atmosphericLight.G, transmission);
float b = RecoverChannel(sourcePixel.B / 255f, atmosphericLight.B, transmission);
\end{lstlisting}

\subsection{Adaptive omega}

\subsubsection{Động cơ sử dụng omega thích nghi}

Trong DCP cổ điển, \(\omega\) thường được chọn là một hằng số, ví dụ \(0.95\). Cách chọn này thuận tiện nhưng khó phù hợp với mọi vùng trong ảnh. Một giá trị \(\omega\) lớn có thể giúp làm rõ cảnh ở vùng có haze dày, nhưng cũng dễ tạo cảm giác quá gắt ở vùng trời sáng hoặc các bề mặt có độ bão hòa thấp. Ngược lại, nếu giảm \(\omega\) để giữ sự tự nhiên cho các vùng nhạy cảm, mức khử haze ở các vùng cần xử lý mạnh lại có thể chưa đủ.

Notebook \texttt{dehaze-omega-regression.ipynb} tiếp cận vấn đề này bằng cách học một ánh xạ đơn giản từ đặc trưng cục bộ của ảnh sang \(\omega(x)\). Mục tiêu không phải là thay DCP bằng một mô hình phức tạp, mà là học một quy luật điều chỉnh hệ số dehaze theo từng vị trí ảnh.

\subsubsection{Đặc trưng đầu vào và nhãn huấn luyện}

Notebook sử dụng ba đặc trưng cục bộ:
\begin{itemize}
\item dark channel của ảnh haze,
\item độ tương phản cục bộ, tính bằng độ lệch chuẩn trên ảnh xám,
\item độ bão hòa màu.
\end{itemize}

Trong notebook, độ tương phản được tính từ phương sai cục bộ của ảnh xám với trọng số độ chói chuẩn \(0.299, 0.587, 0.114\). Độ bão hòa được tính từ chênh lệch giữa kênh lớn nhất và nhỏ nhất, chuẩn hóa theo cường độ cực đại của điểm ảnh. Các công thức này cũng được sử dụng trong phần hiện thực adaptive omega của hệ thống.

Nhãn huấn luyện \(w_{\text{target}}\) không được gán thủ công mà được suy ngược từ cặp ảnh haze--clear trên tập dữ liệu RESIDE OTS. Với atmospheric light \(A\) ước lượng từ ảnh haze, notebook tính:
\begin{equation}
w_{\text{target}} = \mathrm{clamp}\left(\frac{1 - t_{\text{target}}}{\mathrm{DCNH} + \varepsilon}, 0, 1\right),
\end{equation}
trong đó \(\mathrm{DCNH}\) là dark channel của ảnh haze đã chuẩn hóa theo \(A\), còn \(t_{\text{target}}\) được suy ra từ cặp ảnh haze và clear theo mô hình vật lý. Ý tưởng ở đây là tìm ra giá trị \(\omega\) ``phù hợp nhất'' cho từng vị trí ảnh dưới dạng nhãn mềm trong miền \([0,1]\).

\subsubsection{Mô hình hồi quy và công thức logits}

Mô hình học trong notebook là một lớp tích chập \(1\times1\) với ba đầu vào và một đầu ra. Về bản chất, đây là một hồi quy tuyến tính theo pixel, không phải một mô hình học sâu phức tạp. Đầu ra của mô hình là logits, sau đó được đưa qua hàm sigmoid để thu được \(\omega(x)\):
\begin{equation}
\omega_{\mathrm{map}}(x)=\sigma(z(x))=\frac{1}{1+e^{-z(x)}}.
\end{equation}

Sau huấn luyện, notebook in ra quy luật đã học:
\begin{equation}
z(x) = -1.9900\,D(x) + 1.2305\,C(x) - 2.0765\,S(x) + 1.8193,
\end{equation}
trong đó \(D(x)\) là dark channel, \(C(x)\) là contrast và \(S(x)\) là saturation. Từ đó:
\begin{equation}
\omega_{\mathrm{map}}(x)=\sigma\left(-1.9900\,D(x)+1.2305\,C(x)-2.0765\,S(x)+1.8193\right).
\end{equation}

Xét về trực giác, hệ số dương của contrast làm \(\omega\) tăng ở các vùng có cấu trúc rõ hơn, trong khi hệ số âm của saturation và dark channel làm \(\omega\) giảm ở các vùng màu kém bão hòa hoặc có dark channel cao. Điều này phù hợp với mục tiêu giữ xử lý nhẹ hơn ở các vùng sáng, nhạy cảm hoặc dễ gây cảm giác quá tay.

\subsubsection{Triển khai adaptive omega trong hệ thống}

Trong phần hiện thực, công thức trên không được học lại tại thời gian chạy mà được mã hóa trực tiếp thành một phép tính gọn để dựng \texttt{omegaMap}. Sau khi tính dark channel, contrast và saturation, hệ thống tính logits rồi đưa qua sigmoid để thu được \(\omega(x)\) cho từng điểm ảnh. Cách triển khai này giữ chi phí tính toán ở mức thấp, phù hợp với yêu cầu chạy trong plugin.

Đoạn mã sau thể hiện bước suy diễn \texttt{omegaMap}:

\begin{lstlisting}[style=csharpblock,caption={Trích đoạn suy diễn adaptive omega trong plugin}]
float variance = meanGraySq[i] - (meanGray[i] * meanGray[i]);
float contrast = MathF.Sqrt(MathF.Max(0f, variance));

float cMax = MathF.Max(r, MathF.Max(g, b));
float cMin = MathF.Min(r, MathF.Min(g, b));
float saturation = (cMax - cMin) / (cMax + 1e-6f);

float logits =
    -1.9900f * darkChannel[i] +
     1.2305f * contrast -
     2.0765f * saturation +
     1.8193f;
omegaMap![i] = 1.0f / (1.0f + MathF.Exp(-logits));
\end{lstlisting}

Như vậy, adaptive omega trong hệ thống được triển khai dưới dạng một phép ánh xạ theo pixel với chi phí thấp, trong khi phần huấn luyện và khảo sát hệ số được thực hiện riêng trong notebook.

Tuy vậy, nhánh adaptive vẫn có giới hạn rõ ràng. Bản đồ \(\omega(x)\) hiện chỉ được suy ra từ ba đặc trưng thủ công và một mô hình hồi quy rất gọn, nên không thể bao quát mọi tình huống thực tế phức tạp. Các cảnh có nguồn sáng bất thường, phản xạ mạnh, màu sắc quá cực đoan hoặc phân bố haze khác xa dữ liệu đã dùng trong notebook vẫn có thể cho kết quả chưa tối ưu.

\subsection{Kết hợp nhánh cơ sở và nhánh adaptive}

Hệ thống có hai chế độ vận hành. Ở chế độ \textbf{Fixed}, plugin không tạo \texttt{omegaMap}; toàn bộ ảnh dùng chung một \(\omega\) cơ sở bằng \(0.95\), sau đó được nhân với \texttt{strengthFactor}. Ở chế độ \textbf{Adaptive}, plugin xây dựng \texttt{omegaMap(x)} một lần trong bước cache và dùng bản đồ này khi recover từng điểm ảnh.

Điểm chuyển nhánh nằm ngay trong quá trình xây cache. Nếu cờ \texttt{adaptive} tắt, cache chỉ chứa ảnh nguồn, \texttt{transmissionSlope} và atmospheric light. Nếu cờ này bật, cache chứa thêm \texttt{omegaMap}. Nhờ cấu trúc đó, cùng một transmission nền có thể phục vụ cả hai chế độ mà không làm thay đổi phần lõi của bước recovery.

Về mặt ứng dụng, adaptive mode hữu ích nhất ở các vùng có cường độ sáng lớn, độ bão hòa thấp hoặc biến thiên cấu trúc không đồng đều. Đây là các trường hợp mà một \(\omega\) cố định dễ hoặc là xử lý quá mạnh, hoặc là xử lý chưa đủ. Việc đưa \(\omega\) về dạng bản đồ theo pixel giúp plugin linh hoạt hơn mà vẫn giữ được tinh thần của phương pháp DCP.
